\documentclass[runningheads]{llncs}
\usepackage[T1]{fontenc}
\usepackage{graphicx}
\usepackage{indentfirst}

\title{Impacto da Computação Quântica\\na Criptografia Atual - VPNs}
\titlerunning{Computação Quântica na Criptografia Atual - VPNs}

\author{Miguel Grilo\inst{1} \and Tiago Ramalho\inst{2} \and Diogo Belfo\inst{3}\\
58387 | 58514 | 58706\\
\email{\{l58387, l58514, l58706\}@alunos.uevora.pt}}
\authorrunning{M. Grilo \and T. Ramalho \and D. Belfo}
\institute{Escola de Ciências e Tecnologias, Universidade de Évora, Portugal \\
UC: Segurança Informática\\}

\begin{document}

\maketitle              

\begin{abstract}
A rápida evolução da computação quântica e a iminente chegada do ``Q-Day'' representam uma ameaça existencial para a criptografia clássica de chave pública. As Redes Privadas Virtuais (VPNs), infraestruturas fundamentais para a segurança do tráfego corporativo e governamental remoto, encontram-se particularmente vulneráveis, sendo já alvo ativo de campanhas de interceção do tipo \textit{Harvest Now, Decrypt Later} (HNDL). Em resposta a este paradigma, o \textit{National Institute of Standards and Technology} (NIST) padronizou recentemente novos algoritmos de Criptografia Pós-Quântica (PQC), como o ML-KEM e o ML-DSA. Este artigo analisa o estado da arte desta transição e apresenta uma avaliação crítica dos desafios práticos de engenharia subjacentes à sua implementação. Conclui-se que, embora a robustez matemática dos novos algoritmos esteja consolidada, a sua integração em protocolos amplamente utilizados — como o IPsec, TLS 1.3 e WireGuard — introduz obstáculos operacionais severos. O aumento substancial do tamanho das chaves e das assinaturas provoca sobrecargas na latência do \textit{handshake} e aumenta o risco de fragmentação de pacotes e rejeição por \textit{firewalls} ao nível do \textit{Maximum Transmission Unit} (MTU). Perante este cenário, recomenda-se a adoção urgente de arquiteturas de transição baseadas em esquemas híbridos, sustentadas pelo princípio da agilidade criptográfica, como estratégia imperativa para garantir a resiliência a longo prazo durante a migração dos sistemas legados.

\keywords{Computação Quântica \and Criptografia Pós-Quântica \and Redes Privadas Virtuais (VPN) \and Agilidade Criptográfica \and Esquemas Híbridos.}
\end{abstract}

\section{Introdução}

\subsection{O Paradigma Quântico}
A segurança das comunicações digitais contemporâneas baseia-se em problemas matemáticos de elevada complexidade computacional, como a fatorização de números inteiros e o cálculo de logaritmos discretos. 
Contudo, o paradigma da computação quântica ameaça subverter esta fundação tecnológica. 
Ao contrário da computação clássica, que processa informação em bits determinísticos (0 ou 1), a computação quântica utiliza \textit{qubits} (bits quânticos). Tirando partido de propriedades da mecânica quântica — nomeadamente o entrelaçamento e a superposição — um \textit{qubit} não se limita a representar estritamente 0 ou 1, podendo assumir uma combinação de ambos os estados em simultâneo. É esta capacidade que permite aos computadores quânticos realizar cálculos em paralelo a uma escala exponencial.
Há cerca de duas décadas, a comunidade científica provou teoricamente que este paradigma tem o potencial de quebrar quase toda a criptografia de chave pública atualmente implementada \cite{article}.

\subsection{O Conceito de "Q-Day"}
A eventual construção de um Computador Quântico Criptograficamente Relevante (CRQC) marcará a chegada do chamado ``Q-Day'', o momento em que a atual criptografia de chave pública se tornará obsoleta, colocando em risco iminente a infraestrutura crítica mundial. Michele Mosca introduziu um modelo de avaliação de risco cibernético fundamental para compreender a iminência desta ameaça: se o tempo necessário para reequipar a infraestrutura criptográfica global ($y$) somado ao tempo de vida necessário de confidencialidade dos dados ($x$) for superior ao tempo estimado para a disponibilização de um CRQC ($z$), então a segurança da informação já se encontra, no presente, irremediavelmente comprometida ($x + y > z$) \cite{article}.

\subsection{Problema}
Neste cenário de vulnerabilidade global, as Redes Privadas Virtuais (VPNs) emergem como um dos alvos primordiais de cibercrime. Funcionando como as principais guardiãs do tráfego remoto, as VPNs encapsulam e protegem comunicações empresariais e governamentais críticas na sua travessia por redes não confiáveis \cite{10827179}. A dependência de algoritmos assimétricos clássicos para o estabelecimento destes túneis e para a troca inicial de chaves cria um ponto único de falha severo.

A dimensão deste problema é exacerbada por campanhas do tipo \textit{Harvest Now, Decrypt Later} (HNDL), nas quais intervenientes maliciosos intercetam proativamente e armazenam o tráfego encriptado atual das organizações, aproveitando a exposição em diversas camadas da rede de transporte, para o desencriptarem futuramente quando possuírem os recursos quânticos necessários \cite{kundu2026postquantumcryptographicanalysismessage}. Deste modo, protocolos essenciais à operação de túneis VPN — como o IPsec, o TLS e o WireGuard — encontram-se expostos e sob forte risco de depreciação \cite{oterogarcia2025introducingpostquantumalgorithmsopen}. A preservação da integridade do tráfego remoto e a prevenção da quebra retroativa da confidencialidade impõem, por conseguinte, a análise e a adoção urgente de mecanismos de transição baseados em Criptografia Pós-Quântica (PQC) \cite{10827179}.

\section{Ameaças Quânticas à Criptografia Clássica}

\subsection{Criptografia Assimétrica sob Ataque}
A criptografia assimétrica, ou de chave pública, desempenha um papel fulcral no ecossistema das VPNs, sendo responsável pela autenticação e pelo estabelecimento seguro das chaves de sessão. A sua segurança deriva da intratabilidade de problemas matemáticos tradicionais, como a fatorização de grandes números primos (utilizada no RSA) e a resolução de logaritmos discretos, aplicável ao \textit{Elliptic Curve Cryptography} (ECC) e ao \textit{Diffie-Hellman} (DH). O desenvolvimento do Algoritmo de Shor demonstrou que um computador quântico será capaz de resolver estes problemas em tempo polinomial, desconstruindo integralmente as garantias de segurança destes mecanismos \cite{article}. Consequentemente, a base de confiança e autenticação na qual assentam protocolos como o IKE (no IPsec) e o \textit{handshake} do TLS torna-se matematicamente obsoleta perante a ameaça quântica.

\subsection{Criptografia Simétrica e o Algoritmo de Grover}
Enquanto a criptografia assimétrica enfrenta uma quebra estrutural, o impacto na criptografia simétrica — encarregue de encriptar o volume de tráfego de dados no interior do túnel VPN — é qualitativamente diferente. O Algoritmo de Grover atua sobre algoritmos como o \textit{Advanced Encryption Standard} (AES), oferecendo uma aceleração quadrática em ataques de procura exaustiva (\textit{brute-force}) a chaves criptográficas \cite{article}. Na prática, isto significa que o tempo necessário para quebrar uma chave de $N$ bits num computador quântico é equivalente ao tempo necessário para quebrar uma chave de $N/2$ bits num computador clássico. Embora isto reduza drasticamente a segurança efetiva do padrão AES-128 (rebaixando-o para uma resistência equivalente a 64 bits), a duplicação do tamanho das chaves atua como contramedida eficaz. Por esta razão, a utilização de AES-256 mantém um nível robusto de 128 bits de segurança quântica, sendo ainda amplamente considerado seguro e recomendado para adoção nas arquiteturas de rede atuais e futuras.

\subsection{Ameaça \textit{Store Now, Decrypt Later} (SNDL)}
A inevitabilidade prática da construção de computadores quânticos viáveis sublinha a urgência do risco associado à ameaça \textit{Store Now, Decrypt Later} (SNDL) — frequentemente referida na literatura por \textit{Harvest Now, Decrypt Later} (HNDL). Embora a capacidade computacional para quebrar os algoritmos atuais possa estar a anos de distância, atores maliciosos procedem à recolha sistemática de grandes volumes de tráfego de rede encriptado à medida que este atravessa infraestruturas não confiáveis \cite{kundu2026postquantumcryptographicanalysismessage}. A premissa consiste em armazenar dados cujo valor perdura no tempo — tais como segredos militares, dados médicos, propriedade intelectual e comunicações governamentais confidenciais — para os sujeitar a um ataque futuro baseado no algoritmo de Shor --- um algoritmo quântico que fatora inteiros grandes exponencialmente mais rápido que algoritmos clássicos. Dada a elevada esperança de vida de confidencialidade deste tipo de dados, a implementação de medidas de proteção pós-quânticas não pode aguardar pela maturação tecnológica das ameaças, constituindo um imperativo de mitigação a curto prazo \cite{article}.

\section{Vulnerabilidades nos Protocolos de VPN Atuais}

\subsection{Handshake e Troca de Chaves}
O processo de estabelecimento de um túnel VPN seguro inicia-se sempre com uma fase de \textit{handshake}, durante a qual as partes negociam os parâmetros de segurança e derivam uma chave de sessão partilhada. Em protocolos amplamente implementados, como o IPsec (através do protocolo IKEv2) e o OpenVPN (assente no TLS), este processo depende intrinsecamente de algoritmos como o \textit{Diffie-Hellman} ou o seu equivalente em curvas elípticas (ECDH). Esta fase inicial revela-se o ponto mais crítico e exposto da infraestrutura face à ameaça quântica. Uma vez que as mensagens trocadas durante o \textit{handshake} transitam em texto limpo ou sob proteção assimétrica vulnerável, um adversário com capacidades de interceção a nível de rede (camadas L3 a L6) pode capturar passivamente estes pacotes \cite{kundu2026postquantumcryptographicanalysismessage}. Com o surgimento de um computador quântico, o atacante poderá quebrar retroativamente o problema do logaritmo discreto, extrair a chave de sessão simétrica e, consequentemente, decifrar todo o tráfego que fluiu através daquele túnel específico.

\subsection{PKI (Public Key Infrastructure)}
Para além da confidencialidade da troca de chaves, as VPNs dependem da Infraestrutura de Chaves Públicas (PKI) para assegurar a autenticação mútua entre o cliente e o \textit{gateway}. A integridade deste ecossistema baseia-se em assinaturas digitais, tradicionalmente forjadas com RSA ou ECDSA, que atestam a veracidade dos certificados emitidos pelas Autoridades de Certificação (CAs). A rutura destes algoritmos implica o colapso total do modelo de confiança \cite{article}. Sem assinaturas digitais resistentes, as VPNs tornam-se suscetíveis a ataques ativos de \textit{Man-in-the-Middle} (MitM), nos quais um adversário quântico pode forjar um certificado legítimo e fazer-se passar por um \textit{gateway} corporativo autorizado, intercetando e manipulando o tráfego em tempo real sem que os mecanismos clássicos de segurança detetem a intrusão.

\subsection{VPNs de Acesso Remoto vs. Site-to-Site}
A arquitetura de implementação da VPN dita a complexidade e a urgência da mitigação do risco quântico. As VPNs \textit{Site-to-Site}, tipicamente configuradas sobre IPsec para interligar localizações físicas distintas de uma organização (ex: entre uma sucursal e a sede) \cite{oterogarcia2025introducingpostquantumalgorithmsopen}, agregam volumes massivos de tráfego crítico. O comprometimento de um único túnel \textit{Site-to-Site} expõe a totalidade da comunicação inter-filiais. No entanto, a sua atualização para padrões pós-quânticos é logisticamente mais simples, por envolver apenas a reconfiguração de \textit{routers} ou \textit{firewalls} nas extremidades da ligação.

Em contraste, as VPNs de Acesso Remoto (como aquelas que empregam TLS ou protocolos mais modernos como o WireGuard \cite{10827179}) ligam dispositivos individuais à rede corporativa. O impacto do seu comprometimento afeta essencialmente a privacidade do utilizador e as credenciais de acesso, mas a dificuldade de atualização é exponencialmente maior. A transição exige a distribuição de novos clientes de \textit{software} e a injeção de novas bibliotecas criptográficas em milhares de terminais remotos heterogéneos (\textit{laptops}, \textit{smartphones}, dispositivos IoT), enfrentando desafios significativos de fragmentação, compatibilidade e gestão do ciclo de vida dos \textit{endpoints}.

\section{Criptografia Pós-Quântica (PQC) e Normas}

\subsection{O Papel do NIST}
A transição para a era pós-quântica requer o estabelecimento de um novo paradigma de segurança suportado por algoritmos matemáticos exaustivamente testados e normalizados. O \textit{National Institute of Standards and Technology} (NIST) assumiu a liderança global nesta matéria através do seu processo \textit{Post-Quantum Cryptography Standardization}. Recentemente, este esforço plurianual resultou na publicação das normas federais finais --- destacando-se o padrão FIPS 203 - que especificam as famílias de algoritmos matemáticos (como os baseados em reticulados, ou \textit{lattice-based}) concebidos para resistir aos ataques promovidos tanto por computadores clássicos como quânticos \cite{nist2024fips203}. O trabalho de escrutínio promovido pelo NIST dita o padrão da indústria, constituindo o alicerce fundamental para que o setor tecnológico e as infraestruturas de rede iniciem a migração de algoritmos vulneráveis de forma fidedigna e rigorosa \cite{kundu2026postquantumcryptographicanalysismessage}.

\subsection{Mecanismos de Encapsulamento de Chaves (KEM)}
Para suprimir a vulnerabilidade crítica introduzida pela troca de chaves assimétricas baseada em problemas de logaritmos discretos, a infraestrutura contemporânea convergiu para os Mecanismos de Encapsulamento de Chaves (KEM). Em termos conceptuais, um KEM consiste num aglomerado de algoritmos desenhados estritamente para o estabelecimento de um segredo partilhado (\textit{shared secret}) através de um canal público, garantindo resiliência total face à ameaça quântica \cite{nist2024fips203}. O algoritmo escolhido e normalizado como o padrão dominante para esta finalidade foi o \textit{Module-Lattice-Based Key-Encapsulation Mechanism} (ML-KEM), anteriormente submetido como \textit{CRYSTALS-Kyber}. A transposição do ML-KEM para o processo de \textit{handshake} nas VPNs garante que a chave de sessão criptográfica permanece invulnerável, anulando diretamente o perigo associado aos vetores de interceção antecipada e desencriptação futura \cite{10827179,cryptoeprint:2020/379}.

\subsection{Assinaturas Digitais Pós-Quânticas}
A confidencialidade de um túnel VPN carece de valor se a identidade das partes intervenientes não for verificável de forma segura. Com o iminente declínio do RSA e do ECDSA na autenticação dos \textit{gateways} e infraestruturas PKI, a adoção de Assinaturas Digitais Pós-Quânticas afigura-se inegociável. Para este efeito, o NIST definiu o \textit{Module-Lattice-Based Digital Signature Algorithm} (ML-DSA), derivado primordialmente do esquema \textit{CRYSTALS-Dilithium}, como a norma principal \cite{cryptoeprint:2020/071}. O uso de ML-DSA salvaguarda a integridade dos certificados emitidos perante adversários capazes de forjar identidades, mas obriga a uma restruturação protocolar significativa: as assinaturas e chaves públicas baseadas em reticulados requerem dimensões muito superiores às suas contrapartes clássicas, introduzindo novos paradigmas a nível de processamento em terminais remotos e sobrecargas acentuadas na largura de banda de ligação \cite{cryptoeprint:2020/071}.

\subsection{Agilidade Criptográfica}
A complexidade inerente à adoção de criptografia de nova geração determina que esta transição assuma contornos de evolução contínua em detrimento de uma substituição imediata e pontual. Neste aspeto, o paradigma de "agilidade criptográfica" converteu-se num requisito de topo para arquiteturas de comunicações robustas. Uma VPN caracterizada por agilidade criptográfica está desenhada com alicerces modulares, provendo aos seus gestores a capacidade técnica para suportar, adicionar, substituir ou revogar algoritmos criptográficos sem a obrigação de reconstruir a fundação de \textit{hardware} ou reescrever o \textit{software} do zero \cite{kundu2026postquantumcryptographicanalysismessage}. A adoção desta filosofia, materializada recorrentemente através de configurações híbridas que aliam primitivos testados aos novos algoritmos do NIST (como o ML-KEM), revela-se imperativa para acomodar a interoperabilidade de sistemas legados e colmatar o impacto provável na eventual descoberta futura de fragilidades lógicas em algoritmos recentes \cite{oterogarcia2025introducingpostquantumalgorithmsopen}.

\section{Estratégias de Mitigação e Novas Arquiteturas}

\subsection{Esquemas Híbridos (Abordagem Recomendada)}
A transição imediata e exclusiva para novos algoritmos pós-quânticos acarreta o risco inerente à adoção de primitivas matemáticas recentes, que, apesar de teoricamente imunes à computação quântica, podem conter vulnerabilidades estruturais ou lógicas ainda não descobertas perante ataques clássicos. Por este motivo, a comunidade criptográfica e as entidades de normalização recomendam fortemente a adoção de esquemas híbridos como o caminho de migração transitório ótimo \cite{mallick2026studypostquantumstatus}. Esta abordagem consiste em conjugar, no mesmo processo de \textit{handshake} da VPN, a derivação de chaves através de um algoritmo clássico amplamente provado (como o \textit{Elliptic-Curve Diffie-Hellman} através do X25519) com um Mecanismo de Encapsulamento de Chaves pós-quântico (como o ML-KEM) \cite{oterogarcia2025introducingpostquantumalgorithmsopen}. A premissa central de um esquema híbrido é garantir que a segurança da ligação não seja suscetível ao elo mais fraco; mesmo na eventualidade de o algoritmo PQC ser comprometido no futuro, o túnel VPN manterá, no mínimo, a robustez do algoritmo clássico face a computadores convencionais \cite{mallick2026studypostquantumstatus}.

\subsection{Quantum Key Distribution (QKD)}
Em oposição à Criptografia Pós-Quântica (PQC) --- que se baseia em problemas matemáticos complexos resolvidos por computadores tradicionais ---, a Distribuição de Chaves Quânticas (QKD, do inglês \textit{Quantum Key Distribution}) assenta diretamente nas leis fundamentais da mecânica quântica. Utilizando os princípios da incerteza de Heisenberg e o teorema da não-clonagem, a QKD permite que duas partes partilhem uma chave criptográfica com a garantia teórica absoluta de que qualquer interceção no canal ótico perturba irremediavelmente o estado quântico dos fotões, alertando os intervenientes \cite{article}. Apesar de oferecer segurança incondicional pautada pelas leis da física, a sua aplicação no ecossistema atual de VPNs enfrenta limitações pragmáticas substanciais. A QKD requer infraestruturas dispendiosas de \textit{hardware} ótico dedicado, apresenta severas restrições de distância devido à atenuação do sinal (sendo dependente de nós de repetição confidenciais) e não é viável para estabelecer túneis pontuais em clientes remotos (\textit{software VPNs}) ou dispositivos móveis \cite{kundu2026postquantumcryptographicanalysismessage}.

\subsection{Protocolos Modernos}
A eficácia da proteção pós-quântica dita a sua integração nativa nos protocolos subjacentes às VPNs contemporâneas. Um dos avanços mais consistentes na literatura recai sobre o protocolo WireGuard, amplamente reconhecido pela sua leveza e elevada performance. Trabalhos recentes formalizaram propostas como o \textit{PQ-WireGuard}, uma variante que substitui integralmente a fase de \textit{handshake} original por um modelo genérico assente apenas em KEMs, procurando fornecer tanto confidencialidade \textit{forward secrecy} quântica como autenticação robusta no estabelecimento de túneis VPN \cite{cryptoeprint:2020/379,10827179}. Em paralelo, a indústria tem feito progressos assinaláveis na adaptação de \textit{standards} históricos. O protocolo TLS 1.3 (utilizado pelo OpenVPN) e o IKEv2 (fio condutor do IPsec) possuem mecanismos já padronizados ou em fase avançada de normalização (via IETF) para suportar extensões híbridas de troca de chaves. Não obstante o avanço na derivação de chaves híbridas, a literatura destaca que a transição plena da fase de autenticação para as novas assinaturas pós-quânticas nos referidos protocolos ainda constitui um potencial desafio a nível de eficiência de rede e adoção alargada \cite{mallick2026studypostquantumstatus,cryptoeprint:2020/071}.

\section{Impacto Prático e Desafios de Engenharia}

\subsection{Performance e Latência}
A adoção intrínseca de algoritmos pós-quânticos introduz uma sobrecarga (\textit{overhead}) computacional e de rede assinalável. As assinaturas digitais e as chaves públicas baseadas em reticulados, nomeadamente o ML-DSA e o ML-KEM, possuem dimensões consideravelmente superiores às suas contrapartes clássicas (RSA ou ECDSA). Este aumento drástico do tamanho do \textit{payload} reflete-se diretamente no prolongamento do tempo necessário para completar o \textit{handshake} criptográfico em protocolos de túnel essenciais, como o TLS 1.3 e o IKEv2. Adicionalmente, estudos de performance empírica atestam que a latência no estabelecimento das ligações sofre degradações tangíveis, um impacto que é agravado pelo tempo adicional de processamento em CPU para verificação e encapsulamento das chaves complexas \cite{cryptoeprint:2020/071,oterogarcia2025introducingpostquantumalgorithmsopen}. Em infraestruturas VPN com restrições rigorosas de tempo real e disponibilidade, esta sobrecarga pode constituir um impedimento à adoção fluida.

\subsection{Fragmentação de Pacotes}
Um dos desafios técnicos mais críticos a nível de engenharia de redes decorre das limitações relativas ao \textit{Maximum Transmission Unit} (MTU). Em virtude do peso substancial dos certificados e das novas estruturas de chave, as mensagens transitadas excedem frequentemente o limite convencional de um pacote não fragmentado (tipicamente fixado em 1500 bytes para tráfego UDP em protocolos de VPN baseados em túneis). Esta discrepância volumétrica impõe a fragmentação forçada dos pacotes na camada IP para que estes consigam transitar na rede. No entanto, o ecossistema tecnológico contemporâneo revela-se hostil a este comportamento: diversas \textit{firewalls}, \textit{routers} de fronteira e \textit{middleboxes} encontram-se rigorosamente configurados para descartar tráfego IP fragmentado como medida preventiva de segurança. Consequentemente, a fragmentação induzida pela PQC origina perdas de pacotes sucessivas, culminando em falhas abortadas de \textit{handshake} e inviabilizando o estabelecimento do túnel VPN \cite{cryptoeprint:2020/071,mallick2026studypostquantumstatus}.

\subsection{Gestão de Ciclo de Vida}
A resposta à ameaça quântica suplanta largamente a mera atualização de primitivas criptográficas em código-fonte, configurando-se como um imenso desafio logístico na gestão do ciclo de vida da infraestrutura de TI. A obrigatoriedade de migrar sistemas abrange milhares de \textit{endpoints} inerentemente heterogéneos, estendendo-se desde os \textit{gateways} centralizados até a uma vasta dispersão de dispositivos móveis, \textit{laptops} e sistemas IoT limitados em recursos. Asseverar a convivência transitória entre clientes já capacitados para a PQC e sistemas legados exige políticas de interoperabilidade intrincadas e suportadas por agilidade criptográfica. Observando o modelo matemático de Michele Mosca, o tempo físico e o esforço de orquestração despendidos para substituir e propagar estas atualizações universalmente ($y$) materializam o maior calcanhar de Aquiles das organizações. É imperativo encetar e planear este ciclo de renovação atempadamente, salvaguardando a infraestrutura antes que a janela de segurança ($x + y > z$) se feche de modo irreversível \cite{article,kundu2026postquantumcryptographicanalysismessage}.

\section{Conclusão}

\subsection{Balanço entre a ameaça teórica e a prontidão tecnológica}
A análise efetuada ao longo deste artigo demonstra que a ameaça quântica à criptografia clássica transitou de uma mera conjetura teórica para um risco cibernético palpável. O paradigma emergente de campanhas \textit{Harvest Now, Decrypt Later} (HNDL) atesta que o tráfego atual de túneis VPN já se encontra exposto à interceção de agentes estatais, comprometendo de forma severa dados com uma longa esperança de vida de confidencialidade \cite{kundu2026postquantumcryptographicanalysismessage,article}. Em resposta a este problema, a comunidade científica alcançou um marco histórico de prontidão tecnológica, cimentado pelo esforço do NIST na finalização das normas FIPS relativas a mecanismos baseados em reticulados, como o ML-KEM e o ML-DSA \cite{nist2024fips203}. Não obstante a robustez matemática dos novos algoritmos estar consolidada, o balanço global revela que a engenharia de infraestruturas de rede acusa um atraso significativo. A latência induzida e, fundamentalmente, os constrangimentos operacionais decorrentes da fragmentação de pacotes e rejeição por \textit{firewalls} ao nível do MTU evidenciam que a prontidão criptográfica não se traduz, ainda, numa viabilidade sistémica de adoção imediata e indolor \cite{cryptoeprint:2020/071,mallick2026studypostquantumstatus}.

\subsection{Recomendações para organizações: quando começar a migração para VPNs Quantum-Resistant}
Face às evidências, a inércia não constitui uma opção de governação técnica aceitável. Retomando o teorema de avaliação de risco formulado por Michele Mosca ($x + y > z$), o esforço de orquestração e o tempo exigido para a reestruturação e migração de redes corporativas ($y$) prolongam-se frequentemente por anos \cite{article}. Por conseguinte, as organizações devem iniciar de imediato a auditoria às suas infraestruturas criptográficas e delinear o planeamento orçamental para a transição. A recomendação central assenta na adoção ágil e eficiente de esquemas híbridos nos túneis VPN, garantindo desde já mitigação quântica e salvaguardando a resiliência contra vulnerabilidades criptoanalíticas clássicas ainda não descobertas \cite{oterogarcia2025introducingpostquantumalgorithmsopen,mallick2026studypostquantumstatus}. Além disso, é imperativo que futuros \textit{deployments} garantam o princípio primordial da agilidade criptográfica, permitindo trocas modulares nos ecossistemas das VPNs corporativas sem a necessidade de dispendiosas reengenharias. O começo do "Q-Day" adverte o colapso iminente da confiança digital, e a robustez das comunicações de amanhã dependerá estritamente da proatividade estratégica iniciada hoje.

\newpage

\bibliographystyle{splncs04}
\bibliography{bibliography}

\end{document}